
%
\section{}      %
Turing's machines (TMs) %
are one of the most studied models of computation.
More than fifty years ago, Hennie proved quadratic lower bounds \cite{DBLP:journals/iandc/Hennie65}
for solving simple problems on one-tape TMs. While this remains the
best time lower bound for one-tape machines, lower bounds for stronger
models have been established in many papers, %
 including \cite{PPST83,MR808746,MaS87,KalyanasundaramS92,DBLP:conf/coco/Santhanam01,DBLP:journals/tcs/MelkebeekR05,DBLP:journals/cc/Williams06},
some of which are discussed below.

In 1994, Impagliazzo, Nisan, and Wigderson constructed a pseudorandom
generator \cite{INW94} that \emph{fools }one-tape TMs. More precisely
they show how to efficiently map a seed of length $O(\sqrt{t})$ into
a string of length $t$ that cannot be distinguished from uniform
by a one-tape Turing machine running in time $t$.

However,   %
their result does not apply to \emph{randomized} Turing machines,
%
which %
can ``toss coins.'' 
The model in \cite{INW94} corresponds
to randomized machines that begin by filling the tape with coin tosses
in a one-way fashion, and then operate deterministically. By contrast,
randomized machines can toss coins at any time during a computation.
Equivalently, we can think of randomized machines as having an extra
one-way tape with random bits on it. Note that in \cite{INW94} the
Turing machine receives the output of the generator on its (only)
tape. A randomized machine instead receives it on a separate tape.
This allows the machine to perform several tasks \emph{in linear time},
such as copying the random bits on the work tape in reverse order,
comparing the first half of the random bits with the second half,
etc. All of these operations require \emph{quadratic time} on a basic
one-tape machine. And in fact, an instantiation of the \cite{INW94}
generator can be broken by a randomized TM  %
(see \expref{Section}{sec:RTM-breaks-INW}).

In this paper we give the first generator for randomized TMs. Since
there are several variants of TMs let us first define the model precisely
and then state our result.
\begin{defn}
A randomized Turing machine (RTM) is a Turing machine with two tapes:

- A two-way, read-write work tape,

- and a one-way random tape.
\end{defn}

\begin{thm}
\label{thm:prg} There are explicit pseudorandom generators with seed
length $O(\sqrt{t}\log^{3}t\log q)$ and error $\eps=(tq)^{-\sqrt{t}}$
that fool RTMs with $q$ states running in time $t$. That is, no
such RTM can tell whether its random tape is filled with uniform bits
or with the output of the generator on a uniform
random  %
seed, except with advantage $\eps$.
\end{thm}

%

The seed length is, up to the polylogarithmic factors, the best possible
short of proving super-quadratic lower bounds for TMs.

\medskip

One motivation for constructing pseudorandom generators is proving
lower bounds for stronger models. The RTM model above is the natural
randomized extension of the basic one-tape machine model. Now we consider
a different model to extend: Turing machines with a work tape and
a two-way read-only input tape\emph{.} This is one of the strongest
models for which we can prove lower bounds. Time lower bounds of the
form $n^{1+\Omega(1)}$ were shown in \cite{MaS87,DBLP:journals/tcs/MelkebeekR05,DBLP:journals/cc/Williams06}.
The question of extending these lower bounds to hold even for randomized
machines was asked by several authors. In particular it was asked
in \cite{Viola-PhDthesis-journal-hack,ViolaBPvsE} (by the author),
where lower bounds for an intermediate model are proved, and in a
blog post \cite{Lipton-old-ideas-and-lower-bounds-on-sat} by Lipton.

In this paper we prove a lower bound for this randomized model. Again,
let us first define the model.
\begin{defn}
An RTM2 is a Turing machine with three tapes:

- A two-way, read-write work tape,

- a one-way random tape, and

- a two-way, read-only input tape.
\end{defn}

We prove a lower bound for a problem in $\Sigma_{3}\Timee(n)$ -- linear
time with two quantifier alternations. By standard completeness results
the lower bound holds for $QSAT_{3}$ -- the problem of deciding
the validity of a formula with 2 quantifier alternations. This for
example follows from the proof of Theorem 7 in \cite{DvM06}. Previously
lower bounds were not known even for functions computable in exponential
time $E=\Timee(2^{O(n)})$.

%
\begin{thm}
\label{thm:RTM2-lb}There is a function computable in $\Sigma_{3}\Timee(n)$
which requires time $n^{1+\Omega(1)}$ on any RTM2 with bounded error
probability.
\end{thm}

The lower bound holds for general \emph{two-sided} error. Moreover,
as in previous work, see for example \cite{DBLP:journals/tcs/MelkebeekR05},
it in fact holds even if the machine has \emph{direct access}
(a.\,k.\,a.   %
random access) to the input. This means that there is a separate index
tape and in one time step the machine can move the input head to the
position written in binary on the index tape.

\section{Overview of techniques \label{sec:Overview-of-techniques}}

The proof of \expref{Theorem}{thm:prg} relies on a new simulation of
RTMs by a family of space-efficient \emph{branching programs}, henceforth
called \emph{programs} for brevity. It is critical that each program
in the family is allowed to read bits in a different order, so let
us first define such programs.
\begin{defn}
A \emph{(branching) program }with space $s$ on $t$ bits consists
of a permutation $\pi$ of the bits and a layered directed graph with
$t+1$ layers. Each layer has $2^{s}$ nodes, except the first layer
which has $1$, and the last layer which has $2$. Each node in layer
$i\le t$ has two edges, labeled with $0$ and $1$, connecting to
nodes in layer $i+1$. Nodes on layer $m+1$ are labeled with $0$
and $1$. On a $t$-bit input, the program follows the path corresponding
to reading the input in a one-way fashion according to the permutation
$\pi$: layer $i$ reads input bit $\pi(i)$. It then outputs the
label of the last node.
\end{defn}

If $M$ is an RTM we write $M(r)$ for the output of $M$ when its
random tape is initialized with $r$ (and the work tape is blank).
We can now state our simulation.
\begin{thm}
\label{thm:M-to-OR-BP}Let $M$ be an RTM running in time $t$ and
with $q$ states. There is a family of $(tq)^{O(\sqrt{t})}$ programs
$P_{\alpha}$ using space $O(\sqrt{t}\log tq)$ such that for any
$r$:

- if $M(r)=1$ then there is exactly one $\alpha$ such that $P_{\alpha}(r)=1$,

- if $M(r)=0$ then $P_{\alpha}(r)=0$ for every $\alpha$.
\end{thm}

The string $\alpha$ in the simulation is used to guess short \emph{crossing
sequences} which partition the work tape into small blocks; the program
then simulates the machine one block at the time. This idea goes back
at least to the work by Maass and Schorr \cite{MaS87}, see also the
paper by van Melkebeek and Raz \cite{DBLP:journals/tcs/MelkebeekR05}.
We show that it can be applied to RTMs if the random bits can be permuted.
Here is where we use that the programs read bits in any order. This
step is inspired by the works of Maass \cite{MR808746} and Kalyanasundaram
and Schnitger \cite{KalyanasundaramS92} who give explicit functions
(of the random bits) that are hard for RTMs. Our partition of the
work tape has the additional property that it can be ``verified''
by the program, a property which is not apparent in previous simulations.
This allows us to guarantee that $P_{\alpha}(r)=1$ for at most one
$\alpha$, which is used in obtaining pseudorandom generators.

\medskip

Given the simulation, we can use pseudorandom generators for programs
that read bits in any order. However, we need a good dependence on
the number of states and the error. Using a result from \cite{HLV-bipnfp}
(Corollary 20 and surrounding discussion) gives a pseudorandom generator
for RTMs with seed length $t^{5/6+o(1)}$. This is sufficient for
\expref{Theorem}{thm:RTM2-lb}. Using the generator in the follow-up by
Forbes and Kelley \cite{ForbesKelley-2018} we improve $5/6$ to $1/2$.

%
%
%
%
%
%
%
\begin{proof}[Proof of \expref{Theorem}{thm:prg} from
  \expref{Theorem}{thm:M-to-OR-BP}]  %
The generator in \cite{ForbesKelley-2018}, Corollary 4.3, fools branching
programs on $t$ bits using space $s$ with error $\eps$ and seed length
$O(\log(t2^{s}/\eps)\log^{2}t)=O(s+\log(1/\eps))\log^{2}t$. We set $\eps$
to $(tq)^{-O(\sqrt{t})}$ and recall $s=O(\sqrt{t}\log tq)$ to obtain
seed length at most $O(\sqrt{t}\log^{3}t\log q)$. To show correctness,
let $U$ be the uniform distribution and $D$ the output distribution
of the generator. We have
%
%
\[
\left|\E[M(U)]-\E[M(D)]\right|=
\left|\E\left[\sum_{\alpha}P_{\alpha}(U)\right]-
\E\left[\sum_{\alpha}P_{\alpha}(D)\right]\right|\le(tq)^{O(\sqrt{t})}\eps
%
=(tq)^{-\Omega(\sqrt{t})}   %
\]
concluding the proof.
%
%
%
%
\end{proof}
%
The proof of \expref{Theorem}{thm:RTM2-lb} is in \expref{Section}{sec:Proof-of-lower-bound}.
It follows closely the proof of Theorem 6.2 in \cite{ViolaBPvsE}
which establishes a lower bound for a different model of Turing machines.
\expref{Theorem}{thm:RTM2-lb} in this paper is to Theorem 6.2 in \cite{ViolaBPvsE}
as \expref{Theorem}{thm:prg} in this paper is to the generator in \cite{INW94}:
Theorem 6.2 in \cite{ViolaBPvsE} does not allow for a separate random-bit
tape, but instead concerns Turing machines whose work tape is initially
filled with random bits.

In turn, the proof from \cite{ViolaBPvsE} relies on a number of other
results in the literature, which are also required for the proof in
this paper. We use the arguments mentioned earlier in \cite{MaS87,DBLP:journals/tcs/MelkebeekR05}
to simulate TMs using little space and non-determinism. We then follow
the lead of Diehl and van Melkebeek \cite{DvM06} who proved time-space
lower bounds for randomized space-bounded computation. They used Nisan's
pseudorandom generator \cite{Nis92} and the Sipser--Gacs--Lautemann
\cite{Sip83b,Lau83} simulation of BPP in the polynomial-time hierarchy.
We use instead the pseudorandom generator given by \expref{Theorem}{thm:prg}.
Here it is critical that the seed length has a good dependence on
the number of states, since we will hard-wire the input to an RTM2
in the states. Turning to the Sipser--Gacs--Lautemann simulation, we
note that this simulation (which has two quantifiers) is too slow
for our purpose; we use the faster simulation in \cite{ViolaBPvsE}
(with three quantifiers). Finally, we rely on the machinery of time-space
tradeoffs; for a survey of those see van Melkebeek \cite{Melkebeek06}
or Williams' thesis \cite{Williams07}.

\begin{table}
\begin{centering}
\begin{tabular}{c||c|c|c|c|c|c|c|c|c|}
\hline 
Tape cell & \multicolumn{1}{c||}{1} & \multicolumn{1}{c||}{2} & \multicolumn{1}{c||}{3} & \multicolumn{1}{c||}{4} & \multicolumn{1}{c||}{5} & \multicolumn{1}{c||}{6} & \multicolumn{1}{c||}{7} & \multicolumn{1}{c||}{8} & 9\tabularnewline
\hline 
\hline 
 & $\star1$ &  &  &  &  &  &  &  & \tabularnewline
\hline 
 &  & H &  &  &  &  &  &  & \tabularnewline
\hline 
 &  &  & H &  &  &  &  &  & \tabularnewline
\hline 
 &  &  &  & H &  &  &  &  & \tabularnewline
\hline 
 &  &  &  &  & H &  &  &  & \tabularnewline
\hline 
 &  &  &  &  &  & H &  &  & \tabularnewline
\hline 
 &  &  &  &  &  &  & $\star3$ &  & \tabularnewline
\hline 
 &  &  &  &  &  &  & $\star3$ &  & \tabularnewline
\hline 
 &  &  &  &  &  & H &  &  & \tabularnewline
\hline 
 &  &  &  &  &  &  & H &  & \tabularnewline
\hline 
 &  &  &  &  &  & H &  &  & \tabularnewline
\hline 
 &  &  &  &  & H &  &  &  & \tabularnewline
\hline 
 &  &  &  & H &  &  &  &  & \tabularnewline
\hline 
 &  &  & H &  &  &  &  &  & \tabularnewline
\hline 
 &  & H &  &  &  &  &  &  & \tabularnewline
\hline 
 &  &  & H &  &  &  &  &  & \tabularnewline
\hline 
 &  &  &  & H &  &  &  &  & \tabularnewline
\hline 
 &  &  & H &  &  &  &  &  & \tabularnewline
\hline 
 &  & H &  &  &  &  &  &  & \tabularnewline
\hline 
 &  &  & $\star2$ &  &  &  &  &  & \tabularnewline
\hline 
 &  &  &  & H &  &  &  &  & \tabularnewline
\hline 
 &  &  &  &  & $\star3$ &  &  &  & \tabularnewline
\hline 
 &  &  &  &  &  & H &  &  & \tabularnewline
\hline 
 &  &  &  &  &  &  & H &  & \tabularnewline
\hline 
 &  &  &  &  &  &  &  & H & \tabularnewline
\hline 
 &  &  &  &  &  &  & H &  & \tabularnewline
\hline 
 &  &  &  &  &  &  &  & H & \tabularnewline
\hline 
 &  &  &  &  &  &  & H &  & \tabularnewline
\hline 
 &  &  &  &  &  &  &  & H & \tabularnewline
\hline 
 &  &  &  &  &  &  &  &  & $\star3$\tabularnewline
\hline 
\hline 
$\text{block}$ & \multicolumn{1}{c||}{1} & \multicolumn{1}{c||}{2} & \multicolumn{1}{c||}{2} & \multicolumn{1}{c||}{2} & \multicolumn{1}{c||}{3} & \multicolumn{1}{c||}{3} & \multicolumn{1}{c||}{3} & \multicolumn{1}{c||}{3} & 3\tabularnewline
\hline 
\hline 
 & \multicolumn{1}{c||}{$b_{1}$} & \multicolumn{1}{c||}{} & \multicolumn{1}{c||}{} & \multicolumn{1}{c||}{$b_{2}$} & \multicolumn{1}{c||}{} & \multicolumn{1}{c||}{} & \multicolumn{1}{c||}{} & \multicolumn{1}{c||}{} & $b_{3}$\tabularnewline
\hline 
\end{tabular}
\par\end{centering}
\caption{Computation table of an RTM with 9 work tape cells reading $6$ random
bits. Each row corresponds to a different time stamp and shows the
position of the head H on the work tape. The symbol $\star$ indicates
when a random bit is read. We have three boundaries shown at the bottom.
The ``block'' row shows the partition of work cells in blocks. The
induced partition on the random-bit tape is $133233$.\label{tab:Computation-table-of}}
\end{table}


\section{Proof of \expref{Theorem}{thm:M-to-OR-BP} \label{sec:Proof-of-simu}}
%

The reader may want to refer to \expref{Table}{tab:Computation-table-of}
for an example of an RTM computation and some of the parameters in
this proof. For integers $i$ and $j$ we write $[i,j]$ for the set
$\{i,i+1,\ldots,j\}$. Because the machine runs in time $t$ it can
only use $t+1$ work cells. We identify these cells with $[1,t+1]$.
There are $t$ \emph{boundaries} between these cells. Boundary $i\in[t]$
is between work cell $i$ and $i+1$.

\paragraph{The string $\alpha$.}

The string $\alpha$ contains $\sqrt{t}$ boundaries $b_{1},b_{2},\ldots,b_{\sqrt{t}}\in[t]$.
These boundaries partition the work cells into $\sqrt{t}+1$\emph{
blocks}, where block $i$ are the cells $[b_{i-1}+1,b_{i}]$ with
$b_{0}=0$ and $b_{\sqrt{t}+1}=t+1$. Each boundary $b_{i}$ belongs
to the set 
\[
B_{i}:=[\sqrt{t}(i-1)+1,\sqrt{t}i].
\]
This implies that each block has $\le2\sqrt{t}$ cells.

The string $\alpha$ also contains $\sqrt{t}$ \emph{crossings $c_{1},c_{2},\ldots,c_{\sqrt{t}}$.}
A crossing $c$ is a tuple $(i\to j,h,q)$ which means that the machine
is crossing from block $i$ to block $j$ with the head on the random
tape in position $h$ and state $q$. The length of $\alpha$ 
is $\sqrt{t}O(\log tq)$ as required.

\paragraph{The property of $P_{\alpha}$.}

We shall give programs $P_{\alpha}$ such that $P_{\alpha}(r)=1$
if and only if

(1) $M(r)=1$,

(2) for every $i$ the boundary $b_{i}$ is the one among $B_{i}$
such that the computation $M(r)$ crosses it the least number of times,
picking the smallest $b_{i}$ in case of ties, and

(3) the crossings $c_{i}$ are those induced by the $b_{i}$ and the
computation $M(r)$.

\paragraph{Concluding the proof assuming the property.}

Assuming we have $P_{\alpha}$ as above, we conclude the proof as
follows. First, we need to verify that there are boundaries with respect
to which the computation has $\le\sqrt{t}$ crossings. This would
be true even if we picked the $b_{i}$ in a progression $b_{i}=\sqrt{t}(i-1)+j$:
Because different values of $j$ give disjoint crossings, and the
total number of crossings is at most the computation time $t$, there
is a value $j\in[1,\sqrt{t}]$ for which such $b_{i}$ induce $\le\sqrt{t}$
crossings. More so it holds for our definition of $b_{i}$.

Also, there cannot be $\alpha\ne\alpha'$ such that $P_{\alpha}(r)=P_{\alpha'}(r)=1$.
This is true because $\alpha$ is uniquely specified by the computation
$M(r)$.

\paragraph{Designing the $P_{\alpha}$.}

It remains to exhibit the programs $P_{\alpha}$. The crossings $c_{i}=(k_{i}\to k_{i+1},h_{i},q_{i})$
induce a partition of the random tape in at most $\sqrt{t}+1$ intervals.
Interval $i$ is $[h_{i-1}+1,h_{i}]$ corresponding to the computation
in block $k_{i}$ from crossing $i-1$ to $i$, where crossing $0$
is the beginning of the computation and $h_{0}=0$, and $h_{\sqrt{t}+1}=t$.
We use the convention that $[x+1,x]$ is the empty interval, which
may arise if the machine goes from one crossing to the next without
reading random bits. These intervals can have any length.

The program simulates the machine one block at the time, reusing space
across the blocks. The program will read the random bits in the following
order. First it reads all the random bits in the intervals corresponding
to work-tape block $1$, then all the intervals corresponding to work-tape
block 2, and so on.

For each block $i$, the program goes through the crossings involving
block $i$ in order and simulates the machine starting when a crossing
enters block $i$ until it leaves it. While doing so it verifies that
the crossings leaving the block are correct assuming those entering
it are. If not the program aborts and outputs zero. For this the program
just needs memory for the work cells in one block, and the state of
the machine, overall $O(\sqrt{t}+\log q)$ bits. Moreover, the program
computes the number of crossings for each boundary in the block $i$
it simulates. This takes additional $O(\sqrt{t}\log t)$ bits. At
the end of the simulation of one block, the program can use this information
to verify that the boundaries match the intended values. Recall that
block $i$ is $[b_{i-1}+1,b_{i}]$ and that each $b_{i}\in B_{i}$.
The computation on block $i$ can verify ``one side'' of the requirements
for both $b_{i-1}$ and $b_{i}$. Specifically, the program verifies
that for each $b'\in B_{i-1}$ that is bigger than $b_{i-1}$ the
number of crossings at $b'$ is at least as big as the number of crossings
at $b_{i-1}$, and that for each $b'\in B_{i}$ that is smaller than
$b_{i}$ the number of crossings at $b'$ is strictly bigger than
the number of crossings at $b_{i}$. If one of these checks does not
pass the program aborts and outputs zero. If all checks pass for every
block then the $b_{i}$ are correct.

Overall the space required by the program is $O(\sqrt{t}+\log q)+O(\sqrt{t}\log t)=O(\sqrt{t}\log tq)$.

\section{Proof of \expref{Theorem}{thm:RTM2-lb} \label{sec:Proof-of-lower-bound}}
%

In this section we present the proof of \expref{Theorem}{thm:RTM2-lb}.
The proof involves several different computational models, so we need
a definition.
\begin{defn}
We denote by $\RTMTime{t}$ the class of problems that can be solved
by an RTM2 machine running in time $t$.

We denote by $\TiSp{s}{t}$ the class of problems that can be solved
in time $t$ and simultaneously space $s$. The precise machine model
is not important here -- we can assume a RAM machine.

For a complexity class $C$ and a function $b$ we denote by $\exists^{b}C$
the class of problems that can be solved in $C$ with an $\exists$
quantifier on $b$ bits: $\exists^{b}C=\{L':\exists L\in C\text{ such that }x\in L'\Leftrightarrow\exists y\in\zo^{b(|x|)}:(x,y)\in L\}$.
We use the corresponding definition for the $\forall$ and probabilistic
BP quantifiers.
\end{defn}

Towards a contradiction, assume that $\Sigma_{3}\Timex{n}\subseteq\RTMTime{n^{1+o(1)}}$.
Let $m:=n^{2}$. We derive the following contradiction as in \cite{ViolaBPvsE}:
\begin{eqnarray}
\Sigma_{3}\Timex{m} & \subseteq & \RTMTime{m^{1+o(1)}}\label{E:TSWriteProof1}\\
 & \subseteq & \BP{m^{1-\Omega(1)}}\exists^{m^{1-\Omega(1)}}\TiSp{\poly(m)}{m^{1-\Omega(1)}}\label{E:TSWriteProof2}\\
 & \subseteq & \exists^{m^{1-\Omega(1)}}\forall^{m^{1-\Omega(1)}}\exists^{m^{1-\Omega(1)}}\TiSp{\poly(m)}{m^{1-\Omega(1)}}\label{E:TSWriteProof3}\\
 & \subseteq & \Sigma_{O(1)}\Timex{m^{1-\Omega(1)}}\label{E:TSWriteProof4}\\
 & \subseteq & \Sigma_{3}\Timex{o(m)}\label{E:TSWriteProof5}\\
 &  & \mathrm{contradiction.}\label{E:TSWriteProof6}
\end{eqnarray}

The only inclusion that requires a new proof is \eqref{E:TSWriteProof2}.
Before providing that proof we review why the other inclusions hold.

\eqref{E:TSWriteProof1} is by padding.

\eqref{E:TSWriteProof3} is by the time-efficient simulation of probabilistic
time with three alternations, proved in \cite{ViolaBPvsE} (Theorem
1.3).

\eqref{E:TSWriteProof4} follows by the fact, usually attributed to
\cite{Nep70}, that one can trade alternations for time in sublinear-space
computations. The required inclusion can be found in Section 3.2 in
\cite{vMe04}, or in \cite{FLvMV05}.

\eqref{E:TSWriteProof5} follows by noting that the assumption $\Sigma_{3}\Timex{n}\subseteq\RTMTime{n^{1+o(1)}}$
plus the simulation in \cite{ViolaBPvsE} (Theorem 1.3) imply that
$\Sigma_{3}\Timex{n}\subseteq\Pi_{3}\Timex{n^{1+o(1)}}$. This allows
us to efficiently collapse the polynomial hierarchy at the third level.

The contradiction in \eqref{E:TSWriteProof6} is with the time hierarchy,
see for example Section 3.1 in \cite{vMe04}.

\subsection{Proving \eqref{E:TSWriteProof2}}

In this subsection we give the proof of inclusion \eqref{E:TSWriteProof2}.
First we need that the generator is computable in linear space and
polynomial time.
\begin{claim}
\label{claim:G-small-space} Given an index $i$ to an output bit,
and a seed, we can compute the output bit $i$ of the generator in
\expref{Theorem}{thm:prg} in space $t^{1-\Omega(1)}$ and simultaneously
polynomial time.
\end{claim}

\begin{proof}
{[}Sketch{]} The generator in \cite{ForbesKelley-2018}, Corollary
4.3, involves computing small-bias generators \cite{NaN93,AGHP92}
and bounded-independence generators \cite{ChG89,ABI86}. Those generators
are in fact computable even in logarithmic space, see Theorem 14 in
\cite{HeV06}. The rest of the computation amounts to taking bit-wise
ANDs and XORs of strings, which can be done efficiently in small space.
\end{proof}
We can now prove Inclusion \eqref{E:TSWriteProof2}. We obtain the
following variant of Claim 6.3 in \cite{ViolaBPvsE}. For an RTM2
machine $M$ we write $M(x;r)$ for the output of $M$ on input $x$
and random bits $r$. Recall that $m=m(n)=n^{2}$.
\begin{claim}
\label{claim:RTM-in-TiSp} Let $M$ be an RTM2 machine with $O(1)$
states running in time $m^{1+o(1)}$. Let $G:\zo^{m^{1-\Omega(1)}}\to\zo^{m^{1+\Omega(1)}}$
be the generator from \expref{Claim}{claim:G-small-space}. Then the language
$\{(x;\sigma):M(x;G(\sigma))=1\}$ is in $\exists^{m^{1-\Omega(1)}}\TiSp{\poly(m)}{m^{1-\Omega(1)}}$.
In particular,
\[
\RTMTime{m^{1+o(1)}}\subseteq\BP{m^{1-\Omega(1)}}\exists^{m^{1-\Omega(1)}}\TiSp{\poly(m)}{m^{1-\Omega(1)}}.
\]
\end{claim}

\begin{proof}
The ``in particular'' part follows from the first part by the correctness
of the generator, using $m^{1-\Omega(1)}$ uniform bits as seed for
the generator. Here note that we hardwire the input $x$ by multiplying
the number of states by $|x|$, and think of the resulting machine
as an RTM.

To prove the first part, denote by $\tau$ the work tape of $M$.
Let us divide $\tau$ in consecutive blocks where the first block
is of size $d$ and all the others are of size, say, $\sqrt{m}$.
The parameter $d$ will be specified later. Note $d$ completely specifies
this subdivision in blocks. For fixed $d$, similarly to before let
us define a \emph{crossing} to be a tuple 
\[
(a\rightarrow a',hi,hr,q)
\]
where $a$ is an index to a block in $\tau$, $a'=a\pm1$, $hi$ is
the position of the head on the input tape, $hr$ is the position
of the head on the random tape, and $q$ is the state. A crossing
$(a\rightarrow a',hi,hr,q)$ corresponds to $M$ crossing the boundary
of block $a$ towards block $a'$ with state $q$ and input-tape and
random-tape head positions $hi$ and $hr$ (the position of the head
on the work tape is determined by $a$ and $a'$).

Since different values of $d\in\{0,1,\ldots,\sqrt{m}-1\}$ give rise
to disjoint sets of blocks, there is $d<b$ such that the number of
crossings induced by the computation $M(x;G(\sigma))$ is at most
$t/\sqrt{m}\leq m^{0.51}$.

\emph{Simulation:} We simulate the machine $M$ as follows: First
we guess a shift $d$ and a sequence of $t/\sqrt{m}$ crossings. Then
we check consistency of the computation one block at the time. For
every block number $i$, we use $\sqrt{m}$ bits of space to simulate
the machine on that block. First we initialize those bits to $0$.
Then we scan, in order, the list of guessed crossings. Whenever we
encounter a crossing of the form $(a\rightarrow i,hi,hr,q)$ we do
the following. We move the input head to $hi$ and we initialize a
counter $r$ to $hr$. We simulate $M$ starting in state $q$ until
it crosses the block boundary towards block $a'$. Whenever $M$ asks
for a random bit, we reply with the $r$-th output bit of $G(\sigma)$
and increase $r$ by one. Then we look at the next crossing $(\alpha\rightarrow\alpha',hi',hr',q')$
in the guessed list and check that $\alpha=i$, $\alpha'=a'$ and
$r=hr'$ and that $hi'$ and $q'$ match too. We continue in this
way until the list is over. Finally, without loss of generality we
can check that the last crossing has the accepting state of $M$.

\emph{Complexity:} Note that each crossing can be specified by $O(\log m)$
bits, and therefore we guess at most $m^{0.51}O(\log m)=m^{1-\Omega(1)}$
bits total. The rest of the computation can be done simultaneously
in time $\poly(m)$ and space $m^{1-\Omega(1)}$. To see this, note
that the generator is computable in these resources by \expref{Claim}{claim:G-small-space},
while the rest of the simulation only uses space for one block, which
takes $\sqrt{m}$ bits, plus $O(\log m)$ bits for various counters.
\end{proof}

\section{RTM breaks INW \label{sec:RTM-breaks-INW}}

In this section we show that an instantiation of the generator in
\cite{INW94} for Turing machines does not fool RTMs. Let $s$ be
a parameter and $G$ be a graph on vertex set $\zo^{s}$. The basic
step of the \cite{INW94} generator is showing the pseudorandomness
of the distribution $(x,y,z)$, where $(x,z)$ is a uniform edge in
$G$ and $y$ is a uniform string in $\zo^{s}$. The strength of the
generator is related to the expansion of $G$. A simple graph with
nearly maximum expansion is the one where $(x,z)$ is an edge when
the bit-wise XOR $v$ of $x$ and $z$ has zero inner product modulo
$2$, that is $v_{1}v_{2}\oplus v_{3}v_{4}\oplus\cdots\oplus v_{s-1}v_{s}=0$.
This corresponds to the distribution $(x,y,x\oplus v)$ where $x$
and $y$ are uniform and $v$ is a uniform string with zero inner
product modulo $2$. (It can be verified directly that the distribution
$v$ has \emph{bias} $2^{-\Omega(s)}$ in the sense of Naor and Naor
\cite{NaN93}, and it is folklore that this equals the normalized
second largest eigenvalue in $G$, which is a well-known algebraic
measure of expansion, see \cite{HooryLW06}.)

We now claim that an RTM can distinguish this distribution from uniform
in time $O(s)$. The RTM first copies $x$ on the work tape, then
moves the work-tape head back to the beginning. Then by XORing corresponding
random bits and work-tape bits the RTM can recover the bits of $v$
and compute the inner product of $v$. Note we use one-way access
to the random tape, but two-way access to the work tape.

By contrast, a one-tape TM cannot distinguish this distribution from
uniform unless it runs in time $\Omega(s^{2})$ \cite{INW94}.

\paragraph{Acknowledgments.}

We thank the anonymous referees for the detailed and helpful feedback.
